% Options for packages loaded elsewhere
\PassOptionsToPackage{unicode}{hyperref}
\PassOptionsToPackage{hyphens}{url}
\documentclass[
]{article}
\usepackage{xcolor}
\usepackage[margin=1in]{geometry}
\usepackage{amsmath,amssymb}
\setcounter{secnumdepth}{-\maxdimen} % remove section numbering
\usepackage{iftex}
\ifPDFTeX
  \usepackage[T1]{fontenc}
  \usepackage[utf8]{inputenc}
  \usepackage{textcomp} % provide euro and other symbols
\else % if luatex or xetex
  \usepackage{unicode-math} % this also loads fontspec
  \defaultfontfeatures{Scale=MatchLowercase}
  \defaultfontfeatures[\rmfamily]{Ligatures=TeX,Scale=1}
\fi
\usepackage{lmodern}
\ifPDFTeX\else
  % xetex/luatex font selection
\fi
% Use upquote if available, for straight quotes in verbatim environments
\IfFileExists{upquote.sty}{\usepackage{upquote}}{}
\IfFileExists{microtype.sty}{% use microtype if available
  \usepackage[]{microtype}
  \UseMicrotypeSet[protrusion]{basicmath} % disable protrusion for tt fonts
}{}
\makeatletter
\@ifundefined{KOMAClassName}{% if non-KOMA class
  \IfFileExists{parskip.sty}{%
    \usepackage{parskip}
  }{% else
    \setlength{\parindent}{0pt}
    \setlength{\parskip}{6pt plus 2pt minus 1pt}}
}{% if KOMA class
  \KOMAoptions{parskip=half}}
\makeatother
\usepackage{color}
\usepackage{fancyvrb}
\newcommand{\VerbBar}{|}
\newcommand{\VERB}{\Verb[commandchars=\\\{\}]}
\DefineVerbatimEnvironment{Highlighting}{Verbatim}{commandchars=\\\{\}}
% Add ',fontsize=\small' for more characters per line
\usepackage{framed}
\definecolor{shadecolor}{RGB}{248,248,248}
\newenvironment{Shaded}{\begin{snugshade}}{\end{snugshade}}
\newcommand{\AlertTok}[1]{\textcolor[rgb]{0.94,0.16,0.16}{#1}}
\newcommand{\AnnotationTok}[1]{\textcolor[rgb]{0.56,0.35,0.01}{\textbf{\textit{#1}}}}
\newcommand{\AttributeTok}[1]{\textcolor[rgb]{0.13,0.29,0.53}{#1}}
\newcommand{\BaseNTok}[1]{\textcolor[rgb]{0.00,0.00,0.81}{#1}}
\newcommand{\BuiltInTok}[1]{#1}
\newcommand{\CharTok}[1]{\textcolor[rgb]{0.31,0.60,0.02}{#1}}
\newcommand{\CommentTok}[1]{\textcolor[rgb]{0.56,0.35,0.01}{\textit{#1}}}
\newcommand{\CommentVarTok}[1]{\textcolor[rgb]{0.56,0.35,0.01}{\textbf{\textit{#1}}}}
\newcommand{\ConstantTok}[1]{\textcolor[rgb]{0.56,0.35,0.01}{#1}}
\newcommand{\ControlFlowTok}[1]{\textcolor[rgb]{0.13,0.29,0.53}{\textbf{#1}}}
\newcommand{\DataTypeTok}[1]{\textcolor[rgb]{0.13,0.29,0.53}{#1}}
\newcommand{\DecValTok}[1]{\textcolor[rgb]{0.00,0.00,0.81}{#1}}
\newcommand{\DocumentationTok}[1]{\textcolor[rgb]{0.56,0.35,0.01}{\textbf{\textit{#1}}}}
\newcommand{\ErrorTok}[1]{\textcolor[rgb]{0.64,0.00,0.00}{\textbf{#1}}}
\newcommand{\ExtensionTok}[1]{#1}
\newcommand{\FloatTok}[1]{\textcolor[rgb]{0.00,0.00,0.81}{#1}}
\newcommand{\FunctionTok}[1]{\textcolor[rgb]{0.13,0.29,0.53}{\textbf{#1}}}
\newcommand{\ImportTok}[1]{#1}
\newcommand{\InformationTok}[1]{\textcolor[rgb]{0.56,0.35,0.01}{\textbf{\textit{#1}}}}
\newcommand{\KeywordTok}[1]{\textcolor[rgb]{0.13,0.29,0.53}{\textbf{#1}}}
\newcommand{\NormalTok}[1]{#1}
\newcommand{\OperatorTok}[1]{\textcolor[rgb]{0.81,0.36,0.00}{\textbf{#1}}}
\newcommand{\OtherTok}[1]{\textcolor[rgb]{0.56,0.35,0.01}{#1}}
\newcommand{\PreprocessorTok}[1]{\textcolor[rgb]{0.56,0.35,0.01}{\textit{#1}}}
\newcommand{\RegionMarkerTok}[1]{#1}
\newcommand{\SpecialCharTok}[1]{\textcolor[rgb]{0.81,0.36,0.00}{\textbf{#1}}}
\newcommand{\SpecialStringTok}[1]{\textcolor[rgb]{0.31,0.60,0.02}{#1}}
\newcommand{\StringTok}[1]{\textcolor[rgb]{0.31,0.60,0.02}{#1}}
\newcommand{\VariableTok}[1]{\textcolor[rgb]{0.00,0.00,0.00}{#1}}
\newcommand{\VerbatimStringTok}[1]{\textcolor[rgb]{0.31,0.60,0.02}{#1}}
\newcommand{\WarningTok}[1]{\textcolor[rgb]{0.56,0.35,0.01}{\textbf{\textit{#1}}}}
\usepackage{longtable,booktabs,array}
\usepackage{calc} % for calculating minipage widths
% Correct order of tables after \paragraph or \subparagraph
\usepackage{etoolbox}
\makeatletter
\patchcmd\longtable{\par}{\if@noskipsec\mbox{}\fi\par}{}{}
\makeatother
% Allow footnotes in longtable head/foot
\IfFileExists{footnotehyper.sty}{\usepackage{footnotehyper}}{\usepackage{footnote}}
\makesavenoteenv{longtable}
\usepackage{graphicx}
\makeatletter
\newsavebox\pandoc@box
\newcommand*\pandocbounded[1]{% scales image to fit in text height/width
  \sbox\pandoc@box{#1}%
  \Gscale@div\@tempa{\textheight}{\dimexpr\ht\pandoc@box+\dp\pandoc@box\relax}%
  \Gscale@div\@tempb{\linewidth}{\wd\pandoc@box}%
  \ifdim\@tempb\p@<\@tempa\p@\let\@tempa\@tempb\fi% select the smaller of both
  \ifdim\@tempa\p@<\p@\scalebox{\@tempa}{\usebox\pandoc@box}%
  \else\usebox{\pandoc@box}%
  \fi%
}
% Set default figure placement to htbp
\def\fps@figure{htbp}
\makeatother
\setlength{\emergencystretch}{3em} % prevent overfull lines
\providecommand{\tightlist}{%
  \setlength{\itemsep}{0pt}\setlength{\parskip}{0pt}}
\usepackage{bookmark}
\IfFileExists{xurl.sty}{\usepackage{xurl}}{} % add URL line breaks if available
\urlstyle{same}
\hypersetup{
  pdftitle={project-charter},
  pdfauthor={kelompok-5},
  hidelinks,
  pdfcreator={LaTeX via pandoc}}

\title{project-charter}
\author{kelompok-5}
\date{2026-08-20}

\begin{document}
\maketitle

\section{PROJECT CHARTER}\label{project-charter}

\subsection{CAMPUS SERVICE
INFRASTRUCTURE}\label{campus-service-infrastructure}

\subsubsection{Implementasi Infrastruktur Aplikasi Layanan Kampus
Berbasis Web Menggunakan Docker, Nginx Reverse Proxy, Database,
Monitoring, Security, dan
Backup}\label{implementasi-infrastruktur-aplikasi-layanan-kampus-berbasis-web-menggunakan-docker-nginx-reverse-proxy-database-monitoring-security-dan-backup}

\begin{center}\rule{0.5\linewidth}{0.5pt}\end{center}

\subsection{1. Identitas Proyek}\label{identitas-proyek}

\begin{longtable}[]{@{}
  >{\raggedright\arraybackslash}p{(\linewidth - 2\tabcolsep) * \real{0.1145}}
  >{\raggedright\arraybackslash}p{(\linewidth - 2\tabcolsep) * \real{0.8855}}@{}}
\toprule\noalign{}
\begin{minipage}[b]{\linewidth}\raggedright
Komponen
\end{minipage} & \begin{minipage}[b]{\linewidth}\raggedright
Keterangan
\end{minipage} \\
\midrule\noalign{}
\endhead
\bottomrule\noalign{}
\endlastfoot
Nama Proyek & Campus Service Infrastructure \\
Judul Proyek & Implementasi Infrastruktur Aplikasi Layanan Kampus
Berbasis Web Menggunakan Docker, Nginx Reverse Proxy, Database,
Monitoring, Security, dan Backup \\
Mata Kuliah & Administrasi Aplikasi Jaringan \\
Metode Pembelajaran & Project Based Learning (PjBL) \\
Program Studi & Teknologi Rekayasa Perangkat Lunak (TRPL) \\
Semester & \ldots\ldots\ldots\ldots\ldots.. \\
Kelas & \ldots\ldots\ldots\ldots\ldots.. \\
Kelompok & Kelompok \ldots\ldots\ldots\ldots\ldots.. \\
Dosen Pengampu & \ldots\ldots\ldots\ldots\ldots.. \\
Asisten & \ldots\ldots\ldots\ldots\ldots.. \\
Periode Proyek & \ldots\ldots\ldots\ldots\ldots.. \\
Repository & \ldots\ldots\ldots\ldots\ldots.. \\
\end{longtable}

\begin{center}\rule{0.5\linewidth}{0.5pt}\end{center}

\section{2. Latar Belakang}\label{latar-belakang}

Perkembangan aplikasi berbasis web membutuhkan infrastruktur jaringan
dan server yang mampu mendukung proses deployment, komunikasi
antarservice, penyimpanan data, keamanan, monitoring, serta pemeliharaan
sistem.

Dalam praktiknya, sebuah aplikasi web tidak hanya membutuhkan source
code aplikasi, tetapi juga membutuhkan lingkungan server yang
terkonfigurasi dengan baik. Komponen seperti Linux Server, container,
network, reverse proxy, database, storage, security, monitoring, dan
backup menjadi bagian penting dalam pengoperasian aplikasi.

Melalui proyek \textbf{Campus Service Infrastructure}, mahasiswa akan
membangun sebuah infrastruktur aplikasi layanan kampus berbasis web
dengan menggunakan Ubuntu Server dan Docker.

Aplikasi web yang digunakan dibuat secara sederhana dan hanya berfungsi
sebagai media untuk menguji implementasi infrastruktur. Dengan demikian,
fokus utama proyek adalah \textbf{administrasi dan implementasi
infrastruktur aplikasi jaringan}, bukan pengembangan aplikasi web yang
kompleks.

\begin{center}\rule{0.5\linewidth}{0.5pt}\end{center}

\section{3. Deskripsi Proyek}\label{deskripsi-proyek}

Proyek ini merupakan implementasi infrastruktur aplikasi layanan kampus
berbasis web menggunakan pendekatan containerized infrastructure.

Infrastruktur yang dibangun meliputi:

\begin{itemize}
\tightlist
\item
  Ubuntu Server
\item
  Git dan repository management
\item
  Docker
\item
  Docker Compose
\item
  Docker Network
\item
  Web/Application Container
\item
  Database Container
\item
  Nginx Reverse Proxy
\item
  Persistent Storage
\item
  HTTPS
\item
  Firewall
\item
  Monitoring
\item
  Backup dan Restore
\item
  Deployment
\item
  Troubleshooting
\end{itemize}

Aplikasi web yang digunakan memiliki fungsi sederhana, misalnya:

\begin{itemize}
\tightlist
\item
  layanan pemesanan,
\item
  layanan dokumen,
\item
  layanan pengajuan,
\item
  layanan informasi,
\item
  atau layanan kampus sederhana lainnya.
\end{itemize}

Pemilihan fungsi aplikasi dapat berbeda antar kelompok, tetapi
\textbf{arsitektur dan komponen infrastrukturnya tetap mengacu pada
standar proyek}.

\begin{center}\rule{0.5\linewidth}{0.5pt}\end{center}

\section{4. Tujuan Proyek}\label{tujuan-proyek}

\subsection{4.1 Tujuan Umum}\label{tujuan-umum}

Membangun dan mengimplementasikan infrastruktur aplikasi layanan kampus
berbasis web yang terintegrasi, terkontainerisasi, aman, dapat
dimonitor, memiliki mekanisme backup dan restore, serta dapat dilakukan
deployment ulang secara terstruktur.

\subsection{4.2 Tujuan Khusus}\label{tujuan-khusus}

Mahasiswa diharapkan mampu:

\begin{enumerate}
\def\labelenumi{\arabic{enumi}.}
\tightlist
\item
  Menyiapkan Ubuntu Server sebagai lingkungan deployment.
\item
  Menggunakan Git untuk version control dan repository management.
\item
  Mengimplementasikan Docker dan Docker Compose.
\item
  Membuat komunikasi antar-container menggunakan Docker Network.
\item
  Mengimplementasikan Nginx sebagai Reverse Proxy.
\item
  Menghubungkan aplikasi dengan database.
\item
  Mengimplementasikan persistent storage.
\item
  Mengkonfigurasi HTTPS/SSL.
\item
  Mengimplementasikan firewall dasar.
\item
  Melakukan monitoring terhadap service dan server.
\item
  Membuat mekanisme backup dan restore.
\item
  Melakukan deployment aplikasi.
\item
  Melakukan testing dan troubleshooting.
\item
  Mendokumentasikan seluruh proses implementasi.
\item
  Menggunakan Git repository sebagai bukti proses pengembangan proyek.
\end{enumerate}

\begin{center}\rule{0.5\linewidth}{0.5pt}\end{center}

\section{5. Ruang Lingkup Proyek}\label{ruang-lingkup-proyek}

\subsection{5.1 Termasuk dalam Proyek}\label{termasuk-dalam-proyek}

\subsubsection{Infrastruktur Server}\label{infrastruktur-server}

\begin{itemize}
\tightlist
\item
  Instalasi dan konfigurasi Ubuntu Server.
\item
  Konfigurasi IP dan hostname.
\item
  Network connectivity.
\item
  Instalasi Git.
\item
  Konfigurasi environment server.
\end{itemize}

\subsubsection{Container}\label{container}

\begin{itemize}
\tightlist
\item
  Instalasi Docker.
\item
  Docker Image.
\item
  Docker Container.
\item
  Docker Compose.
\item
  Docker Network.
\item
  Docker Volume.
\end{itemize}

\subsubsection{Application Service}\label{application-service}

\begin{itemize}
\tightlist
\item
  Menjalankan aplikasi web sederhana.
\item
  Deployment aplikasi menggunakan container.
\item
  Konfigurasi komunikasi aplikasi dengan database.
\end{itemize}

\subsubsection{Database}\label{database}

\begin{itemize}
\tightlist
\item
  Database server.
\item
  Database schema sederhana.
\item
  Konfigurasi database.
\item
  Persistent data.
\item
  Backup database.
\end{itemize}

\subsubsection{Reverse Proxy}\label{reverse-proxy}

\begin{itemize}
\tightlist
\item
  Instalasi/configuration Nginx.
\item
  Routing request menuju application container.
\item
  Konfigurasi port HTTP/HTTPS.
\end{itemize}

\subsubsection{Security}\label{security}

\begin{itemize}
\tightlist
\item
  Firewall.
\item
  Pembatasan port.
\item
  HTTPS/SSL.
\item
  Pengamanan credential.
\end{itemize}

\subsubsection{Monitoring}\label{monitoring}

\begin{itemize}
\tightlist
\item
  Monitoring status server.
\item
  Monitoring container.
\item
  Monitoring service.
\item
  Pemeriksaan log.
\end{itemize}

\subsubsection{Backup dan Restore}\label{backup-dan-restore}

\begin{itemize}
\tightlist
\item
  Backup database/data.
\item
  Restore database/data.
\item
  Pengujian hasil restore.
\end{itemize}

\subsubsection{Testing}\label{testing}

\begin{itemize}
\tightlist
\item
  Connectivity testing.
\item
  Port testing.
\item
  Service testing.
\item
  Application testing.
\item
  Database testing.
\item
  Reverse proxy testing.
\item
  Security testing.
\item
  Backup/restore testing.
\end{itemize}

\subsubsection{Dokumentasi}\label{dokumentasi}

\begin{itemize}
\tightlist
\item
  Project Charter.
\item
  Arsitektur sistem.
\item
  Network design.
\item
  Deployment guide.
\item
  Testing report.
\item
  Troubleshooting guide.
\item
  Dokumentasi repository.
\end{itemize}

\begin{center}\rule{0.5\linewidth}{0.5pt}\end{center}

\section{6. Di Luar Ruang Lingkup}\label{di-luar-ruang-lingkup}

Proyek tidak berfokus pada:

\begin{itemize}
\tightlist
\item
  Pengembangan aplikasi web yang kompleks.
\item
  UI/UX tingkat lanjut.
\item
  Mobile application.
\item
  Artificial Intelligence.
\item
  Machine Learning.
\item
  Microservices dalam skala enterprise.
\item
  Kubernetes.
\item
  Cloud architecture tingkat lanjut.
\item
  High availability.
\item
  Load balancing skala enterprise.
\end{itemize}

Pengembangan aplikasi dibatasi pada \textbf{aplikasi sederhana yang
memiliki fungsi CRUD dasar} sebagai media pengujian infrastruktur.

\begin{center}\rule{0.5\linewidth}{0.5pt}\end{center}

\section{7. Arsitektur Sistem}\label{arsitektur-sistem}

Arsitektur dasar proyek:

\begin{Shaded}
\begin{Highlighting}[]
\NormalTok{                         CLIENT}
\NormalTok{                           │}
\NormalTok{                           │ HTTP / HTTPS}
\NormalTok{                           ▼}
\NormalTok{                  ┌─────────────────┐}
\NormalTok{                  │  Nginx Proxy    │}
\NormalTok{                  │    :80 / :443   │}
\NormalTok{                  └────────┬────────┘}
\NormalTok{                           │}
\NormalTok{                           ▼}
\NormalTok{                  ┌─────────────────┐}
\NormalTok{                  │   Web / App     │}
\NormalTok{                  │    Container    │}
\NormalTok{                  └────────┬────────┘}
\NormalTok{                           │}
\NormalTok{                           ▼}
\NormalTok{                  ┌─────────────────┐}
\NormalTok{                  │    Database     │}
\NormalTok{                  │    Container    │}
\NormalTok{                  └────────┬────────┘}
\NormalTok{                           │}
\NormalTok{                    Persistent Volume}


\NormalTok{        ┌───────────────────────────────┐}
\NormalTok{        │           Monitoring          │}
\NormalTok{        └───────────────────────────────┘}

\NormalTok{        ┌───────────────────────────────┐}
\NormalTok{        │        Backup \& Restore       │}
\NormalTok{        └───────────────────────────────┘}
\end{Highlighting}
\end{Shaded}

\subsubsection{Alur Utama}\label{alur-utama}

\begin{Shaded}
\begin{Highlighting}[]
\NormalTok{Client}
\NormalTok{   │}
\NormalTok{   ▼}
\NormalTok{Nginx Reverse Proxy}
\NormalTok{   │}
\NormalTok{   ▼}
\NormalTok{Application Container}
\NormalTok{   │}
\NormalTok{   ▼}
\NormalTok{Database Container}
\NormalTok{   │}
\NormalTok{   ▼}
\NormalTok{Persistent Storage}
\end{Highlighting}
\end{Shaded}

Komponen monitoring dan backup berfungsi sebagai komponen pendukung
untuk menjaga reliability dan maintainability sistem.

\begin{center}\rule{0.5\linewidth}{0.5pt}\end{center}

\section{8. Teknologi yang Digunakan}\label{teknologi-yang-digunakan}

\begin{longtable}[]{@{}ll@{}}
\toprule\noalign{}
Komponen & Teknologi \\
\midrule\noalign{}
\endhead
\bottomrule\noalign{}
\endlastfoot
Operating System & Ubuntu Server \\
Version Control & Git \\
Repository & GitHub / GitLab \\
Container & Docker \\
Orchestration & Docker Compose \\
Network & Docker Network \\
Reverse Proxy & Nginx \\
Database & PostgreSQL / MariaDB \\
Application & Web Application sederhana \\
Storage & Docker Volume \\
Security & UFW / Firewall \\
HTTPS & SSL/TLS \\
Monitoring & Monitoring Tools \\
Backup & Shell Script \\
\end{longtable}

\begin{center}\rule{0.5\linewidth}{0.5pt}\end{center}

\section{9. Deliverables}\label{deliverables}

Setiap kelompok wajib menghasilkan:

\begin{longtable}[]{@{}rl@{}}
\toprule\noalign{}
No & Deliverable \\
\midrule\noalign{}
\endhead
\bottomrule\noalign{}
\endlastfoot
1 & Project Charter \\
2 & Git Repository \\
3 & Ubuntu Server Environment \\
4 & Docker Environment \\
5 & Docker Compose Configuration \\
6 & Application Container \\
7 & Database Container \\
8 & Docker Network \\
9 & Nginx Reverse Proxy \\
10 & Persistent Storage \\
11 & HTTPS \\
12 & Firewall Configuration \\
13 & Monitoring \\
14 & Backup Script \\
15 & Restore Script \\
16 & Deployment Script \\
17 & Testing Report \\
18 & Troubleshooting Documentation \\
19 & Final Documentation \\
20 & Demo / Presentation \\
\end{longtable}

\begin{center}\rule{0.5\linewidth}{0.5pt}\end{center}

\section{10. Struktur Repository}\label{struktur-repository}

Repository proyek menggunakan struktur:

\begin{Shaded}
\begin{Highlighting}[]
\NormalTok{campus{-}service{-}infrastructure/}
\NormalTok{│}
\NormalTok{├── README.md}
\NormalTok{├── .gitignore}
\NormalTok{├── .env.example}
\NormalTok{│}
\NormalTok{├── docs/}
\NormalTok{│   ├── project{-}charter.md}
\NormalTok{│   ├── architecture.md}
\NormalTok{│   ├── network{-}design.md}
\NormalTok{│   ├── deployment.md}
\NormalTok{│   ├── testing.md}
\NormalTok{│   ├── troubleshooting.md}
\NormalTok{│   └── screenshots/}
\NormalTok{│}
\NormalTok{├── docker/}
\NormalTok{│   ├── docker{-}compose.yml}
\NormalTok{│   ├── app/}
\NormalTok{│   ├── database/}
\NormalTok{│   └── nginx/}
\NormalTok{│}
\NormalTok{├── app/}
\NormalTok{│}
\NormalTok{├── database/}
\NormalTok{│   ├── schema/}
\NormalTok{│   ├── migration/}
\NormalTok{│   └── seed/}
\NormalTok{│}
\NormalTok{├── config/}
\NormalTok{│   ├── nginx/}
\NormalTok{│   ├── firewall/}
\NormalTok{│   ├── network/}
\NormalTok{│   └── server/}
\NormalTok{│}
\NormalTok{├── scripts/}
\NormalTok{│   ├── backup.sh}
\NormalTok{│   ├── restore.sh}
\NormalTok{│   ├── deploy.sh}
\NormalTok{│   └── health{-}check.sh}
\NormalTok{│}
\NormalTok{├── monitoring/}
\NormalTok{│   ├── README.md}
\NormalTok{│   └── config/}
\NormalTok{│}
\NormalTok{├── tests/}
\NormalTok{│   ├── connectivity/}
\NormalTok{│   ├── service/}
\NormalTok{│   ├── security/}
\NormalTok{│   └── performance/}
\NormalTok{│}
\NormalTok{└── backup/}
\NormalTok{    └── .gitkeep}
\end{Highlighting}
\end{Shaded}

\begin{center}\rule{0.5\linewidth}{0.5pt}\end{center}

\section{11. Git dan Repository
Management}\label{git-dan-repository-management}

Setiap kelompok wajib menggunakan Git untuk mengelola seluruh
konfigurasi dan dokumentasi proyek.

Struktur branch:

\begin{Shaded}
\begin{Highlighting}[]
\NormalTok{main}
\NormalTok{ │}
\NormalTok{ └── develop}
\NormalTok{      │}
\NormalTok{      ├── feature/docker}
\NormalTok{      ├── feature/nginx}
\NormalTok{      ├── feature/database}
\NormalTok{      ├── feature/monitoring}
\NormalTok{      └── feature/backup}
\end{Highlighting}
\end{Shaded}

Contoh pembuatan branch:

\begin{Shaded}
\begin{Highlighting}[]
\FunctionTok{git}\NormalTok{ checkout }\AttributeTok{{-}b}\NormalTok{ feature/docker}
\end{Highlighting}
\end{Shaded}

Commit:

\begin{Shaded}
\begin{Highlighting}[]
\FunctionTok{git}\NormalTok{ add .}
\FunctionTok{git}\NormalTok{ commit }\AttributeTok{{-}m} \StringTok{"feat: implement docker environment"}
\end{Highlighting}
\end{Shaded}

Push:

\begin{Shaded}
\begin{Highlighting}[]
\FunctionTok{git}\NormalTok{ push origin feature/docker}
\end{Highlighting}
\end{Shaded}

Setiap anggota harus memiliki kontribusi Git yang dapat diverifikasi.

\begin{center}\rule{0.5\linewidth}{0.5pt}\end{center}

\section{12. Standar Commit}\label{standar-commit}

\begin{longtable}[]{@{}ll@{}}
\toprule\noalign{}
Prefix & Penggunaan \\
\midrule\noalign{}
\endhead
\bottomrule\noalign{}
\endlastfoot
\texttt{feat:} & Penambahan fitur \\
\texttt{fix:} & Perbaikan \\
\texttt{docs:} & Dokumentasi \\
\texttt{config:} & Konfigurasi \\
\texttt{test:} & Pengujian \\
\texttt{refactor:} & Perapihan \\
\texttt{backup:} & Backup dan restore \\
\texttt{deploy:} & Deployment \\
\end{longtable}

Contoh:

\begin{Shaded}
\begin{Highlighting}[]
\FunctionTok{git}\NormalTok{ commit }\AttributeTok{{-}m} \StringTok{"feat: add nginx reverse proxy"}
\FunctionTok{git}\NormalTok{ commit }\AttributeTok{{-}m} \StringTok{"config: configure docker network"}
\FunctionTok{git}\NormalTok{ commit }\AttributeTok{{-}m} \StringTok{"docs: update deployment guide"}
\FunctionTok{git}\NormalTok{ commit }\AttributeTok{{-}m} \StringTok{"test: add service connectivity test"}
\FunctionTok{git}\NormalTok{ commit }\AttributeTok{{-}m} \StringTok{"backup: add database backup script"}
\end{Highlighting}
\end{Shaded}

\begin{center}\rule{0.5\linewidth}{0.5pt}\end{center}

\section{13. Milestone Proyek}\label{milestone-proyek}

Pelaksanaan proyek mengikuti tahapan pembelajaran mingguan:

\begin{longtable}[]{@{}
  >{\raggedleft\arraybackslash}p{(\linewidth - 6\tabcolsep) * \real{0.0667}}
  >{\raggedright\arraybackslash}p{(\linewidth - 6\tabcolsep) * \real{0.2556}}
  >{\raggedright\arraybackslash}p{(\linewidth - 6\tabcolsep) * \real{0.3667}}
  >{\raggedright\arraybackslash}p{(\linewidth - 6\tabcolsep) * \real{0.3111}}@{}}
\toprule\noalign{}
\begin{minipage}[b]{\linewidth}\raggedleft
Minggu
\end{minipage} & \begin{minipage}[b]{\linewidth}\raggedright
Milestone
\end{minipage} & \begin{minipage}[b]{\linewidth}\raggedright
Target Utama
\end{minipage} & \begin{minipage}[b]{\linewidth}\raggedright
Output
\end{minipage} \\
\midrule\noalign{}
\endhead
\bottomrule\noalign{}
\endlastfoot
1 & Project Initiation & Menentukan proyek dan kebutuhan & Project
Charter + Repository \\
2 & Docker Environment & Menyiapkan Docker & Docker Environment \\
3 & Multi-Container & Menjalankan beberapa service & Docker Compose \\
4 & Network + Reverse Proxy & Menghubungkan service & Docker Network +
Nginx \\
5 & Persistent Storage & Menyimpan data secara persistent & Docker
Volume \\
6 & Database Service & Menyiapkan database & Database Container \\
7 & Security + HTTPS & Mengamankan service & Firewall + HTTPS \\
8 & Monitoring & Memantau service & Monitoring \\
9 & Backup \& Restore & Melindungi data & Backup + Restore \\
10 & Deployment & Menyiapkan deployment & Deployment Script \\
11 & Testing & Menguji sistem & Testing Report \\
12 & Troubleshooting & Menangani masalah & Troubleshooting Report \\
13 & Integration & Mengintegrasikan seluruh komponen & Integrated
System \\
14 & Final Testing & Validasi sistem & Final Test Report \\
15 & Documentation & Finalisasi dokumentasi & Final Documentation \\
16 & Presentation & Demonstrasi proyek & Demo + Presentation \\
\end{longtable}

\begin{center}\rule{0.5\linewidth}{0.5pt}\end{center}

\section{14. Kriteria Keberhasilan}\label{kriteria-keberhasilan}

Proyek dinyatakan berhasil apabila:

\subsubsection{Server}\label{server}

\begin{itemize}
\tightlist
\item
  Ubuntu Server berjalan.
\item
  Server dapat diakses.
\item
  Network connectivity berhasil.
\end{itemize}

\subsubsection{Docker}\label{docker}

\begin{itemize}
\tightlist
\item
  Docker berjalan.
\item
  Docker Compose berjalan.
\item
  Application container berjalan.
\item
  Database container berjalan.
\item
  Container network berjalan.
\end{itemize}

\subsubsection{Application}\label{application}

\begin{itemize}
\tightlist
\item
  Aplikasi dapat diakses.
\item
  Aplikasi dapat berkomunikasi dengan database.
\item
  Fungsi CRUD dasar dapat dijalankan.
\end{itemize}

\subsubsection{Reverse Proxy}\label{reverse-proxy-1}

\begin{itemize}
\tightlist
\item
  Nginx berjalan.
\item
  Request dari client diteruskan ke application container.
\item
  Port HTTP/HTTPS dapat digunakan sesuai konfigurasi.
\end{itemize}

\subsubsection{Storage}\label{storage}

\begin{itemize}
\tightlist
\item
  Persistent volume berjalan.
\item
  Data tetap tersedia setelah container restart.
\end{itemize}

\subsubsection{Security}\label{security-1}

\begin{itemize}
\tightlist
\item
  Firewall aktif.
\item
  Hanya port yang diperlukan yang dibuka.
\item
  HTTPS berhasil dikonfigurasi.
\end{itemize}

\subsubsection{Monitoring}\label{monitoring-1}

\begin{itemize}
\tightlist
\item
  Status container dapat diperiksa.
\item
  Status service dapat diperiksa.
\item
  Log dapat diperiksa.
\end{itemize}

\subsubsection{Backup}\label{backup}

\begin{itemize}
\tightlist
\item
  Backup dapat dilakukan.
\item
  Restore dapat dilakukan.
\item
  Data setelah restore dapat diverifikasi.
\end{itemize}

\subsubsection{Deployment}\label{deployment}

\begin{itemize}
\tightlist
\item
  Sistem dapat dijalankan kembali menggunakan repository.
\item
  Konfigurasi deployment terdokumentasi.
\end{itemize}

\subsubsection{Documentation}\label{documentation}

\begin{itemize}
\tightlist
\item
  Dokumentasi lengkap.
\item
  Repository terstruktur.
\item
  Setiap anggota memiliki kontribusi.
\end{itemize}

\begin{center}\rule{0.5\linewidth}{0.5pt}\end{center}

\section{15. Definition of Done}\label{definition-of-done}

Proyek dianggap \textbf{DONE} apabila seluruh kondisi berikut terpenuhi:

\begin{itemize}
\tightlist
\item[$\square$]
  Ubuntu Server dapat diakses.
\item[$\square$]
  Git repository tersedia.
\item[$\square$]
  Docker berjalan.
\item[$\square$]
  Docker Compose berjalan.
\item[$\square$]
  Application container berjalan.
\item[$\square$]
  Database container berjalan.
\item[$\square$]
  Container network berfungsi.
\item[$\square$]
  Nginx Reverse Proxy berfungsi.
\item[$\square$]
  Persistent storage berfungsi.
\item[$\square$]
  HTTPS dikonfigurasi.
\item[$\square$]
  Firewall dikonfigurasi.
\item[$\square$]
  Monitoring berjalan.
\item[$\square$]
  Backup berhasil.
\item[$\square$]
  Restore berhasil.
\item[$\square$]
  Seluruh service telah diuji.
\item[$\square$]
  Troubleshooting telah dilakukan.
\item[$\square$]
  Dokumentasi lengkap.
\item[$\square$]
  Deployment dapat dilakukan kembali.
\item[$\square$]
  Repository dapat digunakan untuk deployment ulang.
\item[$\square$]
  Setiap anggota memiliki kontribusi Git yang dapat diverifikasi.
\end{itemize}

\begin{center}\rule{0.5\linewidth}{0.5pt}\end{center}

\section{16. Risiko Proyek}\label{risiko-proyek}

\begin{longtable}[]{@{}
  >{\raggedright\arraybackslash}p{(\linewidth - 4\tabcolsep) * \real{0.2796}}
  >{\raggedright\arraybackslash}p{(\linewidth - 4\tabcolsep) * \real{0.3441}}
  >{\raggedright\arraybackslash}p{(\linewidth - 4\tabcolsep) * \real{0.3763}}@{}}
\toprule\noalign{}
\begin{minipage}[b]{\linewidth}\raggedright
Risiko
\end{minipage} & \begin{minipage}[b]{\linewidth}\raggedright
Dampak
\end{minipage} & \begin{minipage}[b]{\linewidth}\raggedright
Mitigasi
\end{minipage} \\
\midrule\noalign{}
\endhead
\bottomrule\noalign{}
\endlastfoot
Docker tidak berjalan & Service tidak dapat dijalankan & Periksa service
dan log Docker \\
Konfigurasi network salah & Container tidak saling terhubung & Periksa
Docker Network \\
Nginx salah konfigurasi & Aplikasi tidak dapat diakses & Validasi
konfigurasi Nginx \\
Database gagal & Aplikasi tidak berfungsi & Periksa database log dan
credential \\
Data hilang & Kehilangan data proyek & Backup berkala \\
HTTPS gagal & Service tidak aman & Validasi SSL/TLS \\
Firewall salah konfigurasi & Service tidak dapat diakses & Periksa rule
UFW \\
Container berhenti & Service tidak tersedia & Monitoring dan health
check \\
Repository tidak teratur & Deployment sulit & Standar struktur dan Git
workflow \\
Credential bocor & Risiko keamanan & Gunakan \texttt{.env} dan
\texttt{.gitignore} \\
\end{longtable}

\begin{center}\rule{0.5\linewidth}{0.5pt}\end{center}

\section{17. Batasan Proyek}\label{batasan-proyek}

Proyek memiliki batasan:

\begin{enumerate}
\def\labelenumi{\arabic{enumi}.}
\tightlist
\item
  Aplikasi web hanya menggunakan fungsi sederhana.
\item
  Pengembangan aplikasi bukan fokus utama.
\item
  Infrastruktur dijalankan pada lingkungan Ubuntu Server.
\item
  Database menggunakan PostgreSQL atau MariaDB.
\item
  Deployment menggunakan Docker dan Docker Compose.
\item
  Monitoring disesuaikan dengan tools yang tersedia.
\item
  HTTPS dapat menggunakan sertifikat self-signed untuk lingkungan
  pembelajaran apabila domain publik tidak tersedia.
\item
  Backup menggunakan mekanisme sederhana berbasis Shell Script.
\item
  Infrastruktur ditujukan untuk lingkungan pembelajaran, bukan
  production enterprise.
\end{enumerate}

\begin{center}\rule{0.5\linewidth}{0.5pt}\end{center}

\section{18. Pembagian Tugas Kelompok}\label{pembagian-tugas-kelompok}

\begin{longtable}[]{@{}
  >{\raggedright\arraybackslash}p{(\linewidth - 2\tabcolsep) * \real{0.4559}}
  >{\raggedright\arraybackslash}p{(\linewidth - 2\tabcolsep) * \real{0.5441}}@{}}
\toprule\noalign{}
\begin{minipage}[b]{\linewidth}\raggedright
Role
\end{minipage} & \begin{minipage}[b]{\linewidth}\raggedright
Tanggung Jawab
\end{minipage} \\
\midrule\noalign{}
\endhead
\bottomrule\noalign{}
\endlastfoot
Project Lead & Koordinasi, timeline, integrasi \\
Infrastructure Engineer & Ubuntu, Docker, server \\
Network / DevOps Engineer & Network, Nginx, deployment \\
Database / Application Engineer & Application dan database \\
Security / Monitoring & Firewall, HTTPS, monitoring \\
Documentation / Testing & Testing, dokumentasi, troubleshooting \\
\end{longtable}

\begin{quote}
Jumlah role dapat disesuaikan dengan jumlah anggota kelompok.
\end{quote}

\begin{center}\rule{0.5\linewidth}{0.5pt}\end{center}

\section{19. Identitas dan Komitmen
Kelompok}\label{identitas-dan-komitmen-kelompok}

\subsubsection{Nama Kelompok}\label{nama-kelompok}

\textbf{Kelompok:}
\ldots\ldots\ldots\ldots\ldots\ldots\ldots\ldots\ldots\ldots\ldots\ldots\ldots\ldots\ldots\ldots\ldots.

\subsubsection{Anggota}\label{anggota}

\begin{longtable}[]{@{}rllll@{}}
\toprule\noalign{}
No & Nama & NIM & Role & Tanda Tangan \\
\midrule\noalign{}
\endhead
\bottomrule\noalign{}
\endlastfoot
1 & & & & \\
2 & & & & \\
3 & & & & \\
4 & & & & \\
\end{longtable}

\subsubsection{Repository}\label{repository}

\textbf{URL Repository:}
\ldots\ldots\ldots\ldots\ldots\ldots\ldots\ldots\ldots\ldots\ldots\ldots\ldots\ldots\ldots\ldots\ldots.

\subsubsection{Server}\label{server-1}

\textbf{Hostname:}
\ldots\ldots\ldots\ldots\ldots\ldots\ldots\ldots\ldots\ldots\ldots\ldots\ldots\ldots\ldots\ldots\ldots.

\textbf{IP Address:}
\ldots\ldots\ldots\ldots\ldots\ldots\ldots\ldots\ldots\ldots\ldots\ldots\ldots\ldots\ldots\ldots\ldots.

\subsubsection{Domain / Host}\label{domain-host}

\textbf{Domain / Hostname:}
\ldots\ldots\ldots\ldots\ldots\ldots\ldots\ldots\ldots\ldots\ldots\ldots\ldots\ldots\ldots\ldots\ldots.

\begin{center}\rule{0.5\linewidth}{0.5pt}\end{center}

\section{20. Persetujuan Project
Charter}\label{persetujuan-project-charter}

Dengan menyetujui Project Charter ini, kelompok berkomitmen untuk
melaksanakan proyek sesuai ruang lingkup, milestone, deliverables, dan
kriteria keberhasilan yang telah ditetapkan.

\begin{longtable}[]{@{}llll@{}}
\toprule\noalign{}
Peran & Nama & Tanggal & Persetujuan \\
\midrule\noalign{}
\endhead
\bottomrule\noalign{}
\endlastfoot
Project Lead & & & \\
Anggota & & & \\
Anggota & & & \\
Anggota & & & \\
Dosen Pengampu & & & \\
\end{longtable}

\begin{center}\rule{0.5\linewidth}{0.5pt}\end{center}

\section{21. Prinsip Utama Proyek}\label{prinsip-utama-proyek}

\begin{quote}
\textbf{Repository bukan hanya tempat menyimpan source code, tetapi
menjadi bukti bahwa mahasiswa mampu membangun, mengonfigurasi, menguji,
mendokumentasikan, dan melakukan deployment infrastruktur aplikasi
secara terstruktur.}
\end{quote}

\begin{quote}
\textbf{Aplikasi web merupakan media pengujian, sedangkan fokus utama
proyek adalah implementasi Infrastruktur Aplikasi Jaringan.}
\end{quote}

\end{document}
